% !TeX root = ../../surface-lagrangian-torus-fibration.tex
% 2026-11-21 / v1.0: Full critical leaves を保つ相対葉族補間．
\section{Saddle を動かす reduced 葉族の相対補間}
\label{ke:relative-reduced-foliations-section}

本節では source の critical moment から離れた自由な normal-circle 領域を扱う．
Saddle の位置を固定する補題\ref{ke:moment-dependent-neck-foliation}とは異なり，
full critical leaf と Reeb graph の型を保ちながら saddle を動かす．
二つの first integral の値，critical values，reduced 面積又は第二作用周期を
一致させることは要求しない．

\subsection{必要な葉族と marking の仮定}

$P$ を連結な compact oriented smooth surface with boundary とし，
各境界成分 $C_a$ に互いに素な固定 collar
$A_a\simeq[0,\epsilon)_r\times S^1$ を取る．
$J$ は実開区間，$\mu_*\in J$ は anchor とする．
本節の smooth simple-saddle family とは，smooth function
\[
 f:J\times P\longrightarrow\RR,\qquad f_\mu=f(\mu,\cdot)
\]
の connected-level-set partition $\mathcal F_\mu$ で，次を満たすものとする．
\begin{enumerate}
\item Interior critical points は全て非退化な index-one critical points である．
  一つの connected critical leaf に含まれる critical point は一つだけである．
  Critical leaf は $\figeight$，その小さい飽和近傍は pair of pants となる．
\item 境界 collar 上の葉は固定された circles $r=\mathrm{constant}$ である．
  従ってそこで $f_\mu=b_a(\mu,r)$，$\partial_rb_a\ne0$ と書けるが，
  $b_a$ 自体が $\mu$ に依存しないとは仮定しない．
\item Boundary components，saddles，およびその間の regular annuli の marking は，
  anchor の marking から smooth continuation したものとする．
  任意に別の Reeb-graph permutation を parameter ごとに指定する意味ではない．
\end{enumerate}
Boundary の角を用いる場合は，固定した toric collar 内で先に丸めてから適用する．
必要な操作の支持はその小さい collar の外に取る．

異なる connected critical leaves の critical values が等しくてもよい．
全 critical values の相異性は仮定しない．
ただし同じ connected leaf に二つの saddle がある場合には，その saddle connection が
摂動で解け得るため，以下の定理の対象ではない．
一つの固定 $J$ 全体での無条件の trivialization も主張せず，
常に anchor の近くへ区間を縮めて用いる．

必要に応じて有限個の fixed marked points $q_1,\ldots,q_N\in\operatorname{int}P$ を
追加できる．その場合は，anchor でこれらが互いに異なる regular leaves 上にあり，
boundary collars と critical leaves から離れていると仮定する．
区間を縮めれば，各点の regularity とこの分離を保てる．
この条件は本研究の distinct regular source markings に適合する．

\subsection{一族の全葉を同時に同定する vector field}

\begin{lemma}[Boundary-relative な simple-saddle 葉族の局所自明化]
\label{ke:simple-saddle-family-trivialization}
上の smooth family に対し，$\mu_*$ を含む小さい $J_0\Subset J$ と
smooth diffeomorphisms $\Phi_\mu:P\to P$ が存在し，
\[
 \Phi_{\mu_*}=\id,\qquad
 \Phi_\mu(\mathcal F_{\mu_*})=\mathcal F_\mu
\]
を満たす．$\Phi_\mu$ は小さい全 boundary collars 上で恒等であり，
上記の fixed marked points を指定した場合はそれらも点ごとに固定する．
Full critical leaves，saddle labels，boundary labels，および全 regular annuli の
continuation marking を保存する．各固定 $k$ について
\[
 \|\Phi_\mu-\id\|_{C^k(P)}\le C_k|\mu-\mu_*|
\]
となる．定数は固定 family と anchor に依存してよく，動く stage に一様とはしない．
\end{lemma}

\begin{proof}
Parameter を移して $\mu_*=0$ とする．以下の区間は必要に応じて縮める．
Critical points は非退化なので，implicit function theorem により
互いに交わらない smooth paths $p_v(\mu)$ をなし，
\[
 c_v(\mu)=f_\mu(p_v(\mu)),\qquad
 \dot c_v(\mu)=\partial_\mu f_\mu(p_v(\mu))
\]
である．Compactness と $df_0\ne0$ on the complement of small critical neighborhoods により，
これら以外の critical points は生じない．

まず anchor の各 full critical leaf $L_v$ を，互いに素な小さい飽和領域 $Q_v$ に入れる．
具体的には $f_0^{-1}([c_v(0)-\eta_v,c_v(0)+\eta_v])$ の，$L_v$ を含む成分を取る．
十分小さい $\eta_v>0$ に対して，$Q_v$ は一つの saddle だけを含み，
全 boundary components は regular level circles である．
この領域の存在には，$L_v$ が一つの connected critical leaf であることを使う．
もし全ての $\eta_v\downarrow0$ でその成分が $L_v$ の指定近傍から出るなら，
nested compact connected components の共通部分が，$f_0^{-1}(c_v(0))$ の中で
$L_v$ より大きい connected set になるからである．
異なる $L_v$ の critical values が同じでも，それらの成分は別々に選べる．

$Q_v$ の各 boundary circle に implicit function theorem を適用し，
level $f_\mu=c_v(\mu)\pm\eta_v$ の小さい normal graph に置き換える．
これにより smooth domains $Q_{v,\mu}$ を得る．
Boundary は level circles，内部の critical point は $p_v(\mu)$ のみなので，
内部の各 level の connected component はその boundary を横切らない．
Boundary と同じ level に属する connected leaf も，regular collar の局所積構造により
当該 boundary circle 自体であり，その外へ別の枝が付くことはない．
従って $Q_{v,\mu}$ は必要な範囲で全葉の和である．
これらの pair-of-pants neighborhoods の外には regular annuli だけが残る．
実際，regular component 上の $\nabla f_\mu/|\nabla f_\mu|^2$ の流れを用いれば，
その component は circle と interval の積になる．
従って以下で構成する全体の vector field は，各 annulus の両端も同時に同定する．

次に $P$ 上の smooth basic function $\zeta_\mu$ を作る．
ここで basic は，各 connected leaf 上で定数という意味である．
$Q_{v,\mu}$ 内で
\[
 \zeta_{v,\mu}
 =\dot c_v(\mu)\,\chi_v(f_\mu-c_v(\mu))
\]
とする．$\chi_v$ は $0$ の近傍で $1$，
$(-\eta_v,\eta_v)$ 内に支持された固定 cutoff とする．
Boundary の近くで零なので，$Q_{v,\mu}$ の外へ零で延長しても smooth である．
これは full critical leaf の近傍全体で $\dot c_v(\mu)$ に等しい．

境界 collar では，より小さい collar で $1$ となる radial cutoff $\beta_a(r)$ を用い，
\[
 \zeta_{a,\mu}=\beta_a(r)\,\partial_\mu b_a(\mu,r)
\]
とし，その外へ零で延長する．従って小さい collar 上では
$\zeta_\mu=\partial_\mu f_\mu$ となる．
これは collar 上の first integral の\emph{表示値}を固定せず，葉だけを固定するために必要である．

Fixed marked points を指定した場合は，それぞれの full regular leaf に
互いに素な飽和 annular neighborhoods $A_{i,\mu}$ を取る．
$d_i(\mu)=f_\mu(q_i)$ とし，その内部で
\[
 \zeta_{i,\mu}=\dot d_i(\mu)\,\chi_i(f_\mu-d_i(\mu))
\]
を加える．$\chi_i=1$ near $0$ とし，支持を当該 annulus 内に取る．
これらの supports は critical neighborhoods と boundary collars から離して取れる．
その結果 $\zeta_\mu(q_i)=\partial_\mu f_\mu(q_i)$ となる．

有限個の上の functions の和を $\zeta_\mu$ とする．
その構成から，全 $P$ 上の smooth function $a_\mu$ があり
\begin{equation}\label{ke:eq:basic-critical-velocity}
 d\zeta_\mu=a_\mu\,df_\mu
\end{equation}
となる．Critical points の近傍では $\zeta_\mu$ は定数なので $a_\mu=0$ と取れる．
Boundary collar では
\[
 a_\mu=
 \frac{\partial_r(\beta_a\partial_\mu b_a)}{\partial_rb_a}
\]
であり，分母は零でない．その他の各 support でも，$f_\mu$ の一変数関数を微分すればよい．
Supports の外では $a_\mu=0$ とする．

今度は
\begin{equation}\label{ke:eq:foliation-transport-equation}
 \partial_\mu f_\mu+df_\mu(Y_\mu)=\zeta_\mu
\end{equation}
を満たす smooth vector field を構成する．
各 $p_v(\mu)$ の近くで parameter-dependent Morse coordinates を取り，
\[
 f_\mu\circ\kappa_{v,\mu}=c_v(\mu)+x^2-y^2
\]
とする．$Y_{v,\mu}\circ\kappa_{v,\mu}=\partial_\mu\kappa_{v,\mu}$ と置けば，
$\zeta_\mu=\dot c_v(\mu)$ の小さい近傍で
\eqref{ke:eq:foliation-transport-equation}を満たす．
Parameter-dependent Morse coordinates の存在は，critical point を先に移した上で
Taylor の積分剰余を
$f_\mu-c_v(\mu)=\tfrac12 x^{\mathsf T}H_\mu(x)x$
と書き，signature が一定の symmetric matrix $H_\mu(x)$ を smooth congruence で
$\operatorname{diag}(1,-1)$ に直すことから従う．
Congruence は，$A\mapsto A^{\mathsf T}\operatorname{diag}(1,-1)A$ の微分が
symmetric matrices へ全射であることに implicit function theorem を適用して取れる．
従ってここで新しい parameter-trivialization theorem を外部入力にしてはいない．

Critical neighborhoods の外では固定補助計量の gradient を使い，
\[
 Y^{\rm reg}_\mu
   =\frac{\zeta_\mu-\partial_\mu f_\mu}{|\nabla f_\mu|^2}\nabla f_\mu
\]
とする．Saddle charts と regular region の partition of unity でこれらを貼り合わせる．
式\eqref{ke:eq:foliation-transport-equation}は $Y_\mu$ に関して線型なので，
各局所解を貼り合わせても同じ式を保つ．
Critical points の近傍では regular formula を使わず，Morse chart の解だけを用いる．
小さい boundary collars では numerator が零なので $Y_\mu=0$ であり，
fixed marked points においても $Y_\mu(q_i)=0$ となる．

この parameter-dependent vector field の flow を $\Phi_\mu$ とする．
Compactness と boundary 近傍での消滅により，小さい共通区間で全 $P$ 上に存在する．
\eqref{ke:eq:foliation-transport-equation}と\eqref{ke:eq:basic-critical-velocity}から
\[
 \frac{d}{d\mu}\Phi_\mu^*(df_\mu)
   =\Phi_\mu^*(d\zeta_\mu)
   =(a_\mu\circ\Phi_\mu)\,\Phi_\mu^*(df_\mu)
\]
であり，従って
\begin{equation}\label{ke:eq:foliation-covector-transport}
 \Phi_\mu^*(df_\mu)
 =\exp\!\left(\int_0^\mu a_s\circ\Phi_s\,ds\right)df_0.
\end{equation}
係数は全点で正である．この式は critical points でも成立し，
regular tangent distributions と critical sets を同時に同定する．
Critical leaf の各 smooth arc とその唯一の vertex を合わせると，
connected critical leaf 全体も一つの connected leaf に写る．
逆 flow にも同じことを適用すれば，像が単なる leaf の真部分集合ではなく
全 connected leaf であることが分かる．
Boundary，saddles，regular annuli の marking は identity からの smooth continuation で保たれる．
最後の $C^k$ 評価は $\Phi_0=\id$ と，固定 compact 上での flow の smooth dependence による．
\end{proof}

\subsection{二つの葉族の比較と moment cutoff}

\begin{theorem}[Full critical leaves を保つ相対葉族補間]
\label{ke:saddle-moving-relative-interpolation}
同じ $P$ と boundary collars 上に，二つの smooth simple-saddle families
$\mathcal F^0_\mu,\mathcal F^1_\mu$ を与える．
全 boundary collars では両者の leaf partition が同じ固定 radial partition であり，
anchor では $\mathcal F^0_{\mu_*}=\mathcal F^1_{\mu_*}$ が
\emph{同じ marked partition} として成り立つとする．
各族は上で定めた smooth first integral を持つが，その値が相互に等しい必要はない．
Fixed marked points を保存したい場合には，上記の distinct-regular-leaf 条件を課す．

十分小さい $\mu_*$ の近傍 $J_0$ 上で，boundary-relative な小さい diffeomorphism 族
$h_\mu$ があり，
\[
 h_{\mu_*}=\id,\qquad
 h_\mu(\mathcal F^0_\mu)=\mathcal F^1_\mu.
\]
さらに $\mu_*$ の近傍で $1$，$J_0$ の外で $0$ の smooth moment cutoff $\chi$ に対して，
\[
 \mathcal F^{\rm new}_\mu=\mathcal F^1_\mu
 \quad\text{near }\mu_*,\qquad
 \mathcal F^{\rm new}_\mu=\mathcal F^0_\mu
 \quad\text{outside }J_0
\]
となる smooth simple-saddle family を構成できる．
Boundary/corner の全葉，full critical leaves の連結性と型，
および指定した fixed marked points を保持する．
Saddle の位置と critical values は動いてよい．
\end{theorem}

\begin{proof}
前補題を二つの family に別々に適用し，trivializations
$\Phi^0_\mu,\Phi^1_\mu$ を得る．
Anchor の partitions が同じなので
\[
 h_\mu=\Phi^1_\mu\circ(\Phi^0_\mu)^{-1}
\]
が所要の比較となる．それぞれの first integral の anchor での表示値も，
boundary での表示値も揃える必要はない．
$h_\mu$ は境界と指定点を固定し，$\mu\to\mu_*$ で全固定階に恒等写像へ収束する．

固定補助計量を用いて $h_\mu(p)=\exp_p v_\mu(p)$ と書き，
$h^r_\mu(p)=\exp_p(rv_\mu(p))$，$0\le r\le1$ とする．
区間を縮めれば全 $h^r_\mu$ が恒等写像に $C^1$ で十分近く，
boundary を固定する diffeomorphisms となる．
実際，local diffeomorphism であり，compactness と boundary-relative な
identity homotopy から degree-one covering なので global diffeomorphism である．
固定点では $v_\mu=0$ なので isotopy 全体もその点を固定する．

\[
 \mathcal F^{\rm new}_\mu
   =h^{\chi(\mu)}_\mu(\mathcal F^0_\mu)
\]
と置く．First integral として
$f^{\rm new}_\mu=f^0_\mu\circ(h^{\chi(\mu)}_\mu)^{-1}$ を使える．
各 moment で一つの diffeomorphism の像なので，新しい critical points，
critical-leaf の分裂，regular circles の切断，又は Reeb-graph の組合せ変化は起こらない．
Reeb graph の parameterization は相互に異なってよいが，marked quotient と
fiber partition は指定どおりに一致する．
\end{proof}

\paragraph{Marking の仮定が必要な理由．}
同じ connected critical leaf に複数 saddle がある場合，small perturbation は
その critical values を分離して別々の critical leaves にし得る．
また固定した二点が anchor では同じ leaf 上にあり，別 parameter では異なる leaves 上に
移る family は，その二点を固定する diffeomorphism では自明化できない．
本定理はこれらの変化を許す意味での単なる ``Morse family'' を対象にしない．
他方，互いに無関係な critical leaves の数値としての critical values の一致は障害でない．
証明中の basic drift $\zeta_\mu$ を各 connected critical neighborhood で別々に選んだためである．

点固定の障害は annulus $P=[-1,1]_r\times S^1_\theta$ 上ですでに現れる．
$\beta(r)=1$ near $0$，$\beta=0$ near $\partial P$ とし，
\[
 f_\mu(r,\theta)=r+\mu\beta(r)\cos\theta
\]
と置く．十分小さい $\mu$ では critical point はなく，boundary 葉も固定される．
しかし $q_1=(0,0)$，$q_2=(0,\pi)$ は $\mu=0$ で同じ葉にあり，
$\mu\ne0$ では $f_\mu(q_1)=\mu$，$f_\mu(q_2)=-\mu$ により異なる葉にある．
従って両点を固定する葉族の自明化は存在しない．
これは moment-dependent な Lagrangian 葉族そのものの不存在ではない．

\subsection{Hamiltonian circle reduction への持ち上げ}

\begin{corollary}[自由 normal band 上の actual Lagrangian 補間]
\label{ke:actual-saddle-moving-neck-interpolation}
$(E^4,\omega)$ に自由な Hamiltonian $S^1$ 作用を与え，
$\iota_V\omega=-d\mu$ とする．$J_0$ の閉包は source の critical moments を避けるとする．
その商の対象領域に固定した smooth trivialization
\[
 q:E\longrightarrow E/S^1\simeq J_0\times P
\]
があり，boundary collars は既存の共通 toric germs と一致すると仮定する．
各 fiber に用いる二葉族が前定理の条件を満たせば，
\[
 q^{-1}(\{\mu\}\times L),\qquad L\in\mathcal F^{\rm new}_\mu
\]
は actual $\omega$ に対する Lagrangian torus fibration を与える．
Critical leaves の逆像は $S^1\times\figeight$ で，その smooth strata は isotropic である．
Normal circle とその向き，boundary/corner の全 torus fibers を保ち，
cutoff の支持外では既存 fibration と一致する．
従って支持を非零 moment の小区間に取れば，central source の pinched torus germs は変更しない．
\end{corollary}

\begin{proof}
Regular leaf の逆像は principal $S^1$ bundle over $S^1$ なので $T^2$ である．
Full critical leaf は一つの saddle を持つ $\figeight$ であり，その上の principal circle
bundle も trivial なので，逆像は $S^1\times\figeight$ である．
これらは connected fibers である．

各二次元 smooth stratum の tangent space は $V$ と，
reduced 一次元葉の tangent vector の lift $W$ で張られる．
$d\mu(W)=0$ より
\[
 \omega(V,W)=-d\mu(W)=0.
\]
従って regular torus は Lagrangian であり，singular fiber の二次元 strata も isotropic である．
Critical orbit は一次元なので同じ結論を持つ．
この計算は，葉族が $\mu$ に依存していても各 fixed moment 上で成立する．
同値に，$f^{\rm new}(\mu,p)$ を $E$ へ引き戻すと $S^1$-invariant なので，
その函数は $\mu$ と Poisson commute する．

Normal-cycle marking を明示するには，$h^r_\mu$ の generating vector field を
principal bundle の connection で horizontal lift する．
Moment を保つ equivariant isotopy が得られ，boundary と cutoff 支持外では恒等に取れる．
これは symplectic isotopy であるとは主張しない．
Lagrangian 性は上の actual-form の計算で別に示したものであり，
reduced diffeomorphism の面積保存又は作用周期保存を必要としない．
全葉の逆像を用いるため，saddle fiber を部分的に切ることもない．
\end{proof}

\subsection{Liouville link の葉族への適用と moment の再 parameterization}

Actual な invariant primitive $\lambda$ と
$\iota_Z\omega=\lambda$ を取り，$[Z,V]=0$ とする．
同じ period と向きの normal generator を使えば，
$M=\lambda(V)$ と元の normal moment は
\[
 M=\mu-\rho
\]
という定数差を持つ．$ZM=M$ とし，anchor を $\mu=\mu_*$，
$M_*:=\mu_*-\rho\ne0$ とする．
Liouville time $s$ と実際の normal moment は
\[
 M(s)=e^sM_*,\qquad
 \mu(s)=\rho+e^s(\mu_*-\rho),\qquad
 s(\mu)=\log\frac{\mu-\rho}{\mu_*-\rho}
\]
で対応する．右辺は anchor の小さい近傍で smooth である．
Period 又は向きを変えて比較する場合は，対応する affine moment change を先に行う．
Liouville time を元の $\mu$ と同じ parameter として扱ってはいけない．

Anchor link の reduced 葉を既存の $\mathcal F^0_{\mu_*}$ そのものとする．
Liouville flow は normal action と可換なので，$s(\mu)$ を用いて
reduced surfaces の smooth diffeomorphism 族を誘導し，
それによって $\mathcal F^1_\mu$ を得る．
対象の compact reduced domains と全 boundary collars を含む共通 buffer 上に
flow が定義されていることを適用条件とする．
Corner では両者が同じ full $T^2$ orbit partition であることを別に確認する．
その下で前定理と前系により，source moments を避けた小区間で
Liouville 葉族から既存 neck 葉族へ相対補間できる．

本節で解いたのはこの葉族補間である．
閉じた link 全体の存在・被覆，actual generators の微分評価，
新しい component cut の全 fiber 飽和，および全曲面の距離比較は，
ここでの smooth-family 入力から自動的には従わない．
